\chapter{Implementation}\label{ch:implementation}

This chapter will introduce an implementation of the algorithm presented in
\Cref{sec:algorithm} and explain the design decisions made.

\section{Non-atomic Operations}

Since \cite{Duplicacy} does not make any explicit claims about the atomicity of
the operations clients perform on the repository it must be assumed they are atomic.
This section will examine the feasibility of these assumptions when using a POSIX~\cite{Posix}
file system.

\subsection{Creating chunks and manifest files}

When creating chunks multiple clients may create the same chunk at the same time.
To deal with this situation the repository must resolve the arising copies by
keeping on of them~\cite[Theorem 1]{Duplicacy}.
Furthermore, other clients may access a chunk or manifest file while it is in the
process of being created, which could result in partial and therefore invalid
data being read.

Based on the structure of the repository described in \cite[section `Naming Chunks']{Duplicacy}
it can be assumed that the chunks and manifest files exist as individual files.
To prevent clients from accessing invalid chunks and manifest files they can be first
created in different directories from the ones containing the valid chunks and manifest
files existing on the same file system.
When a client has finished writing a file it can use the `rename'~\cite{Posix}
function available on all POSIX conforming systems to atomically move it to the
other directories.
Another benefit of this approach is that any existing chunks are replaced,
resolving any arising copies by freeing their occupied storage space once any
processes having the file open during the replacement have closed
it+\cite[function `rename', section `DESCRIPTION']{Posix}.
To prevent clients from opening the same file in the directory for temporary chunks
it can be exploited that the file name can be different from the name after
the rename operation, for example by incorporating the name of the client into the file name

\subsection{Renaming chunks and fossils}

Lock-Free Deduplication requires clients being able to rename chunks to fossils
during the fossil collection step and doing the reverse during the fossil deletion
step~\cite[section `Two Step Fossil Deletion']{Duplicacy}.

Besides the chunks existing as individual files it can also be assumed that
chunks and fossils are located on the same file system~\cite[section `Naming Chunks']{Duplicacy}.
Under these conditions the operation can be performed by using the aforementioned
`rename'+\cite{Posix}.
The case where the renaming operation fails because the fossil or chunk was already
renamed by another client~\cite[Theorem 2]{Duplicacy} can be detected by checking for the
`ENOENT' error code~\cite[function `rename', section `ERRORS']{Posix}.

\subsection{Deleting fossils}

In addition to the aforementioned renaming operation the fossil deletion step
also requires clients being able to delete fossils.
Since the fossils exist as individual files they can be deleted by using the
`unlink'+\cite{Posix} function which also available on all POSIX conforming systems.
The case where the deletion operation fails because the fossil was deleted (or renamed)
by another client~\cite[Theorem 2]{Duplicacy} can be detected by checking for the
`ENOENT' error code similar to the `rename' function~\cite[function `unlink', section `ERRORS']{Posix}.

\subsection{Downloading all manifest files}

Both the fossil collection and the fossil deletion steps begin by downloading all
manifest files from the repository~\cite[Algorithm 1 and 2]{Duplicacy}.
This operation requires enumerating the entries present in the directory containing
the manifest files and reading them and the chunks containing the rest of the
chunk references as outlined in \Cref{sec:algorithm}.
Since the enumeration of a directory requires calling the `readdir' function other
clients can asynchronously create or delete files in the directory, which may not
be reflected in the entries returned by it~\cite[function `readdir', section `DESCRIPTION']{Posix}.
While \Cref{thm:singleclient} prevents manifest files from being deleted and chunks
or fossils from being renamed while a client is performing the fossil collection and
deletion steps other clients may create backups and therefore manifest files during
the enumeration of the directory.

\begin{thm}\label{thm:nonatomic-listing}
    The guarantees provided by an atomic directory enumeration operation to the
    Lock-Free Deduplication algorithm can also be provided by using an enumeration
    operation which may not reflect manifest file creations happening during its execution.
\end{thm}

\paragraph{Proof idea} Should the enumeration of all manifest files be an atomic
operation its result would be a snapshot of the repository at a point during
its execution.
The directory enumeration operation possibly not reflecting manifest file creations
happening during its execution would be equivalent to them moved freely before or
after that point from the perspective of the client performing the directory
enumeration operation.
\Cref{fig:listing-proof-idea} contains an example with the snapshot point being marked
as a dotted red line.

\begin{figure}[H]
    \includesvg[width=0.65\textwidth, pretex=\small]{retrieval_cases.drawio.svg}
    \caption{Example of directory operations being moved before or after the snapshot point of an atomic enumeration operation.}\label{fig:listing-proof-idea}
\end{figure}

\paragraph{Proof} Since manifest file creation is not synchronized with the
enumeration operation of the directory an atomic operation can not guarantee
whether a creation operation happening during its execution will be performed
before or after its snapshot point and therefore reflected by its result.
Since a non-atomic enumeration operation provides the same guarantees for
the visibility of individual operations happening during its execution there
is no difference between in this regard.

The same is true for creation operations happening before or after the execution
of the atomic or non-atomic enumeration operations since the non-atomic operation
can only cause creation operations happening during its execution to move before
or after the snapshot point of the atomic operation.

\Cref{fig:manifest-ordering} shows that \nameref{para:policy3} also depends
on the ordering of manifest file creation operations when proceeding with
the fossil deletion step.
Should in the example only the second manifest file creation operation be reflected
by the enumeration operation \nameref{para:policy3} may allow the client performing
the fossil deletion step to continue.
Since \nameref{para:policy3} only ensures that new manifest files created by the
clients can not reference any existing fossils the client has to take all existing
manifest files which where not seen by the fossil collection step into account when
determining whether a fossil should be recovered or deleted.
Since the creation of the first manifest file may not be reflected by the directory
enumeration operation this requirement is violated.

\begin{figure}
    \includesvg[width=0.65\textwidth, pretex=\small]{manifest_ordering.drawio.svg}
    \caption{Example of Policy 3 depending on manifest creation ordering.}\label{fig:manifest-ordering}
\end{figure}

To prevent this situation from occurring when using a non-atomic directory enumeration
operation the client performing the fossil deletion step should perform a second enumeration
operation after \nameref{para:policy3} is satisfied.
Proceeding only after \nameref{para:policy3} is satisfied ensures that manifest
files created after the first enumeration operation completed can not reference
any existing fossils.
This in turn ensures that all manifest files which still reference them where
created before the execution of the second enumeration operation and are therefore
reflected by its result.

\subsection{Accessing a chunk}

Since \nameref{para:policy2} causes chunk access operations to consist of up to
two file access operations another client performing the fossil collection or
deletion steps can rename the chunk or its fossils between them.
This in turn can cause the chunk to appear as missing as shown in \Cref{fig:nonatomic-access}.
In this example a client creates a manifest file called `\#1' during a fossil
collection step, causing on of the chunks referenced by it to be turned into a
fossil to be recovered by the following fossil deletion step.
Any client trying to access the chunk before it is recovered will fail and
access the associated fossil as by \nameref{para:policy2}.
Should the fossil be recovered before the second access operation takes place
it will cause it to fail, making the chunk appear to be missing from the
perspective of the accessing client.

\begin{figure}
    \includesvg[width=0.75\textwidth, pretex=\small]{nonatomic_access.drawio.svg}
    \caption{Example of a chunk appearing as missing.}\label{fig:nonatomic-access}
\end{figure}

\begin{samepage}
    \begin{thm}\label{thm:nonatomic-access}
        Assuming the client does not create any manifest files while accessing a
        chunk or its associated fossil.
        By performing the access pattern in \nameref{para:policy2} two times a client
        will either succeed in accessing either the chunk or its fossil or detect that
        it was deleted.
    \end{thm}
\end{samepage}

\paragraph{Proof} A client performing the fossil deletion step must make sure
\nameref{para:policy3} is satisfied before continuing, causing fossil deletion
steps started after the client created its last manifest file to be delayed until
the client creates a new manifest file.
Since \Cref{thm:singleclient} limits the amount of pending fossil deletion steps
to one the chunk may only be recovered once during the access operations performed
by the client.

Should the client fail to access either the chunk or its associated fossil on
the first try because the fossil was recovered it will succeed in accessing the
chunk on the second try because it can not be recovered a second time.
Should the fossil be deleted instead the client will fail to access it or its
associated chunk on the second try unless it was recreated in the meantime. \\

Another deviation from an atomic access operation is caused by the fact that
after successfully accessing a file with the `open' function it
has to read it using the `read' function, allowing other clients to rename or
delete the file in the meantime~\cite{Posix}.
Since the renaming and deletion operations do not interfere with already opened
files~\cite[section `DESCRIPTION' of the functions `unlink' and `rename']{Posix}
this situation requires no further mitigations.

\section{Error Recovery}

Under real world conditions the algorithm presented in \Cref{sec:algorithm}
may fail at any time, for example in case of a client experiencing a power failure.
This section describes methods to recover from such events.

\subsection{Effects of failure}

The set of operations which modify the repository consists of backup creation
and the fossil collection and deletion steps.

Should a client creating a backup be interrupted it might have already uploaded
some chunks to the repository without creating a manifest file.
Should no client create a manifest file referencing these chunks they will be
invisible to the fossil collection step and therefore remain in the repository
indefinitely.

Similarly, a fossil collection step failing or an existing fossil collection being abandoned
(for example by the client storing it failing to successfully perform the fossil deletion step)
would cause created fossils to remain in the repository.
Should another fossil collection step not try to turn their associated chunks into
fossils again and subsequently remove them during the fossil deletion step they will also
remain in the repository indefinitely.

\subsection{Handling unreferenced chunks}

The unreferenced chunks created by an interrupted backup creation operation can
not simply be deleted because of the possibility of clients creating backups in
parallel and therefore have to be made visible to a fossil collection step.

\begin{thm}\label{thm:unreferenced-chunks}
    Assuming there is a way of enumerating all chunks existing in a repository.
    It is possible to handle unreferenced chunks using the fossil collection and
    deletion steps without disrupting the Lock-Free Deduplication algorithm.
\end{thm}

\paragraph{Proof} Assuming a manifest file referencing all chunks existing in
the repository was created, passing it to the fossil collection step to be deleted
would cause all chunks not referenced by other manifest files to be deleted by the
associated fossil deletion step.

Instead of creating this manifest file the client performing the fossil collection
step can just pretend it already exists and use the set of chunks existing within
the repository as the chunks referenced by it when performing the fossil collection
step, not deleting it afterwards.
Since \Cref{thm:singleclient} prevents other fossil collection and deletion steps
from running in parallel there is no difference between this imaginary manifest file
being deleted by the fossil collection step or it never existing in the first place. \\

A drawback of this solution is that chunks created by backup creation operations
running in parallel will be detected as being unreferenced and therefore turned
into fossils.
While the fossil deletion step will turn them back into chunks the increased
number of renaming operations makes fossil collection steps of this kind more
expensive in comparison to regular fossil collection steps.
A solution to this problem is to only perform the expensive variant on demand
or at a longer interval than the regular fossil collection steps, which is
similar to the recommendations for the solution used by Duplicacy suffering
from the same problem~\cite[topic `Prune command details']{Duplicacy-Forum}.

\subsection{Handling abandoned fossils}

Since abandoned fossils existing within the repository have no associated fossil
collection the timestamp for validating \nameref{para:policy3} can not be determined.

\begin{thm}\label{thm:abandoned-fossils}
    Assuming there is a way of enumerating all fossils existing in a repository.
    It is possible to handle abandoned fossils using the fossil collection and
    deletion steps without disrupting the Lock-Free Deduplication algorithm.
\end{thm}

\paragraph{Proof} A possible solution would be to extend the list of created
fossils recorded by a fossil collection with the fossils already existing in
the repository.
Since they already existed in the repository at the time of the fossil collection
step its associated fossil deletion step will only operate on fossils created
before its fossil collection step completed, which was proven not to disrupt the
Lock-Free Deduplication algorithm by the proof of \Cref{thm:singleclient}.
The fossil deletion step will either recover or delete them, making them no longer
exist in the repository.

To prevent the fossil deletion step from interfering with the proof of 
\Cref{thm:nonatomic-access} by recovering a fossil multiple times the list
of fossils recorded by the fossil collection has to be deduplicated.
